\documentclass[11pt,a4paper]{article}
\usepackage[utf8]{inputenc}
\usepackage{amsmath,amssymb,amsthm}
\usepackage{natbib}
\usepackage{geometry}
\geometry{margin=1in}
\newtheorem{proposition}{Proposition}
\newtheorem{corollary}{Corollary}
\newtheorem{lemma}{Lemma}
\theoremstyle{remark}
\newtheorem{remark}{Remark}
\title{\textbf{Joy of Winning and Loss Aversion in Induced-Value Auctions:\\
A One-Parameter Model of Overbidding in Private and Common Value Settings}}
\author{}
\date{}
\begin{document}
\maketitle
\begin{abstract}
\noindent
Laboratory subjects overbid relative to the risk-neutral Bayesian--Nash benchmark
in private- and common-value auctions. Expectations-based reference dependence of
the K\H{o}szegi--Rabin type is often proposed as an explanation, but its canonical
sealed-bid formulation separates a \emph{good} dimension from a \emph{money}
dimension. In induced-value experiments the object's worth is itself a monetary
redemption value, so the two dimensions collapse. We therefore study
K\H{o}szegi--Rabin preferences with a \emph{single monetary dimension} and a
self-consistent reference point, and add one reference-independent
parameter---a joy-of-winning bonus $w\ge0$ received upon winning. The model nests
pure loss aversion at $w=0$ and yields closed-form equilibria. In the private-value
first-price auction, loss aversion alone already generates overbidding for every
valuation (a first-order risk-aversion channel), at a rate increasing in the
loss-aversion index $\Lambda$; the joy-of-winning bonus amplifies it and produces a
strictly positive bid at a zero valuation, $w/(1-\Lambda)$. In the first-price
wallet game, loss aversion alone has \emph{no} effect on the equilibrium bid
($\beta^*(x)=x$ at $w=0$), while any $w>0$ yields overbidding
$\beta^*(x)-x=w/(1-\Lambda+\Lambda x)$. The contrast is sharp and testable: the
presence of common-value overbidding \emph{beyond} the winner's curse identifies a
joy-of-winning motive, whereas private-value overbidding is consistent with pure
loss aversion. The framework provides a tractable baseline for structural
estimation and richer behavioral auction models.
\end{abstract}


\section{Core Framework: K\H{o}szegi \& Rabin (2006)}

\subsection{Key Innovation and Mathematical Setup}
The reference point is determined endogenously by the agent's \textbf{recent rational expectations} (probabilistic beliefs) rather than an exogenous status quo. Evaluations occur dimension-by-dimension across $K$ separable consumption domains.

For consumption bundle $c = (c_1, \dots, c_K)$ and reference bundle $r = (r_1, \dots, r_K)$:
\begin{equation}
u(c \mid r) = \sum_{k=1}^K m_k(c_k) + \sum_{k=1}^K \mu\big(m_k(c_k) - m_k(r_k)\big),
\end{equation}
where $m_k(\cdot)$ denotes intrinsic consumption utility, and $\mu(\cdot)$ is the piecewise-linear gain-loss utility:
\begin{equation}
\mu(x) = \begin{cases} \eta x & \text{if } x \ge 0, \\ \eta \lambda x & \text{if } x < 0, \end{cases} \qquad (\lambda > 1, \; \eta > 0).
\end{equation}
When outcomes and reference points are distributions $F$ and $G$:
\begin{equation}
U(F \mid G) = \iint u(c \mid r) \, dG(r) \, dF(c).
\end{equation}

\subsection{Solution Concepts}
\begin{enumerate}
    \item \textbf{Personal Equilibrium (PE):} An action plan $d$ is credible if, given that the agent expected to choose $d$, following through is optimal:
    \begin{equation}
    d \in \arg\max_{d'} U\big(F(d') \mid F(d)\big).
    \end{equation}
    \item \textbf{Preferred Personal Equilibrium (PPE):} When multiple actions qualify as PEs, the agent plans the one maximizing ex-ante expected utility:
    \begin{equation}
    d^* = \arg\max_{d \in \text{PE}} U\big(F(d) \mid F(d)\big).
    \end{equation}
\end{enumerate}

\section{Benchmark: The Shoe Shopping Decision}

A consumer decides whether to purchase an item with value $v$ at fixed price $p$. Dimensions: Item ($k=1$) and Money ($k=2$). Set $\eta = 1$, $\lambda > 1$.

\subsection{Equilibrium Analysis}
\begin{itemize}
    \item \textbf{Case 1: Expecting to Buy ($r = (1, -p)$):}
    \begin{align*}
        U(\text{Buy} \mid \text{Expect Buy}) &= v - p, \\
        U(\text{Don't Buy} \mid \text{Expect Buy}) &= -\lambda v + p.
    \end{align*}
    Buying is a PE iff:
    \begin{equation}
    v - p \ge -\lambda v + p \iff p \le \frac{1 + \lambda}{2}v.
    \end{equation}
    The consumer pays above $v$ to avoid the attachment loss of leaving empty-handed.

    \item \textbf{Case 2: Expecting Not to Buy ($r = (0, 0)$):}
    \begin{align*}
        U(\text{Don't Buy} \mid \text{Expect Not}) &= 0, \\
        U(\text{Buy} \mid \text{Expect Not}) &= 2v - (1 + \lambda)p.
    \end{align*}
    Not buying is a PE iff:
    \begin{equation}
    0 \ge 2v - (1 + \lambda)p \iff p \ge \frac{2}{1 + \lambda}v.
    \end{equation}
\end{itemize}

\subsection{Equilibrium Regimes}
\begin{itemize}
    \item $p > \frac{1+\lambda}{2}v$: Unique PE is \emph{Don't Buy}.
    \item $p < \frac{2}{1+\lambda}v$: Unique PE is \emph{Buy}.
    \item $\frac{2}{1+\lambda}v \le p \le \frac{1+\lambda}{2}v$: Multiple equilibria exist. Evaluating ex-ante payoffs selects \textbf{Buy} as the unique PPE iff $p \le v$.
\end{itemize}

\section{Private Value Auction Toy Model}

\subsection{Environment}
Two symmetric bidders $i \in \{1, 2\}$ participate in a First-Price Auction (FPA). Valuations are i.i.d.\ $v_i \sim U[0, 1]$. Dimensions are good ($k=1$) and money ($k=2$). Let $\Lambda \equiv \eta(\lambda - 1) > 0$.

\subsection{Payoff Derivation under PPE}
With planned bid $b$ yielding winning probability $P(b)$:
\begin{align*}
u(c_W \mid c_W) &= v - b \quad (\text{prob } P^2), \\
u(c_L \mid c_L) &= 0 \quad (\text{prob } (1-P)^2), \\
u(c_W \mid c_L) &= (v - b) + \eta v - \eta \lambda b \quad (\text{prob } P(1-P)), \\
u(c_L \mid c_W) &= -\eta \lambda v + \eta b \quad (\text{prob } (1-P)P).
\end{align*}
Aggregating yields the objective function:
\begin{equation}
U(b; v) = \underbrace{P(b)(v - b)}_{\text{Standard Payoff}} - \underbrace{\Lambda P(b)\big(1 - P(b)\big)(v + b)}_{\text{Loss-Aversion Penalty}}.
\end{equation}

\subsection{Equilibrium Strategy}
\begin{proposition}
In the unique symmetric strictly increasing PPE equilibrium, the bidding strategy is:
\begin{equation}
\beta^*(v) = \frac{\frac{1 - \Lambda}{2}v + \frac{2\Lambda}{3}v^2}{1 + \Lambda(1 - v)}.
\end{equation}
\end{proposition}

\begin{proof}
Substituting $P(b) = \beta^{-1}(b)$, the FOC at $b = \beta(v)$ yields the linear ODE:
\begin{equation}
\frac{d}{dv}\Big[\beta(v) \cdot v\big(1 + \Lambda(1 - v)\big)\Big] = (1 - \Lambda)v + 2\Lambda v^2.
\end{equation}
Integrating from $0$ to $v$ with boundary condition $\beta(0) = 0$ gives the result.
\end{proof}

\subsection{Properties}
\begin{itemize}
    \item \textbf{Risk-Neutral Benchmark ($\Lambda = 0$):} $\beta^*(v) = \frac{v}{2}$.
    \item \textbf{High Valuations ($v \to 1$):} $\beta^*(1) = \frac{1}{2} + \frac{\Lambda}{6} > \frac{1}{2}$. Fear of item attachment loss induces \textbf{overbidding}.
    \item \textbf{Low Valuations ($v \to 0$):} $\beta^*(v) \approx \frac{1-\Lambda}{2(1+\Lambda)}v < \frac{v}{2}$ (for $\Lambda < 1$). Fear of monetary expenditure induces \textbf{underbidding}.
    \item \textbf{Second-Price Auction (SPA) Contrast:} SPA introduces uncertain payments conditional on winning. Following Lange and Ratan (2010), truth-telling ($b=v$) fails due to within-dimension price variance.
\end{itemize}

\section{Extension: The Wallet Game (Common Value Auctions)}

\subsection{Baseline Environment}
Two bidders receive private signals $X_i \sim F[0, \bar{x}]$ i.i.d. The common value is $V = X_1 + X_2$. Bidders compete in an English / Second-Price auction where winning payoff at price $P$ is $\pi_i = X_1 + X_2 - P$. Standard symmetric BNE without loss aversion is $b(X_i) = 2X_i$.

\subsection{Incorporating Reference-Dependent Preferences}
\begin{enumerate}
    \item \textbf{Net-Surplus Reference Point ($r = 0$):}
    \begin{equation}
    u(\pi) = \begin{cases} \pi & \text{if } \pi \ge 0, \\ \lambda \pi & \text{if } \pi < 0 \end{cases} \qquad (\lambda > 1).
    \end{equation}
    Losses occur when $X_1 + X_2 < P$ (the winner's curse realization). The penalty slope $\lambda > 1$ amplifies downside risk conditional on winning, depressing equilibrium bids relative to the risk-neutral benchmark.
    
    \item \textbf{Expectations-Based Multi-Dimensional Formulation:}
    \begin{equation}
    u(V, P \mid r_v, r_p) = V - P + \mu_v(V - r_v) + \mu_p(-P - (-r_p)),
    \end{equation}
    where $(r_v, r_p)$ is the rational-expectation joint distribution of value and price induced by symmetric equilibrium strategies.


\end{enumerate}

\section{Extension: The First-Price Wallet Game with Endogenous Reference Dependence}
\label{sec:wallet}
\subsection{Environment}
Two symmetric bidders $i\in\{1,2\}$ compete in a sealed-bid, first-price auction.
Signals are private and i.i.d.,
\begin{equation}
X_i \sim U[0,1], \qquad X_1 \perp X_2,
\end{equation}
and the object has a \emph{common} value
\begin{equation}
V = X_1 + X_2 \in [0,2].
\end{equation}
A bidder who observes $x$, bids $b$, and wins obtains $V-b$; a loser obtains $0$.
As in Section~3, consumption has two separable dimensions: the \emph{good} ($k=1$) and \emph{money} ($k=2$), with gain--loss function $\mu(\cdot)$ and $\Lambda \equiv \eta(\lambda-1)>0$.
\subsection{Risk-Neutral Benchmark}
Suppose the opponent uses the strictly increasing strategy $\beta_0(\cdot)$. A bidder with signal $x$ who bids $b$ wins with probability $B=b^{-1}(b)$ and, conditional on winning, the opponent's signal is $Y\sim U[0,B]$. Expected profit is
\begin{equation}
\int_0^{B}\!\big(x+y-b\big)\,dy = B\Big(x+\tfrac{B}{2}-b\Big).
\end{equation}
The first-order condition at $B=x$ gives $2x-\beta_0(x)-x\beta_0'(x)=0$, i.e.\ $\frac{d}{dx}[x\beta_0(x)]=2x$. With $\beta_0(0)=0$,
\begin{equation}
\boxed{\ \beta_0(x)=x\ } \qquad \text{(first-price wallet game, risk neutral).}
\end{equation}
\subsection{Reference Point and Utility}
We retain the convention of Section~3: the reference point is the consumption
distribution generated by the agent's own planned bid (self-consistent / PPE).
A bidder with signal $x$ who plans a bid $b=\beta(B)$ (which wins with probability $B$) faces
\begin{equation}
\text{Win (prob.\ $B$): good}=x+Y,\ \text{money}=-b; \qquad
\text{Lose (prob.\ $1-B$): good}=0,\ \text{money}=0,
\end{equation}
where $Y\sim U[0,B]$ conditional on winning; the reference is an independent draw from the same distribution. Utility is
\begin{equation}
U = \mathbb{E}\big[m_1+m_2\big] + \mathbb{E}\big[\mu(\Delta_1)\big] + \mathbb{E}\big[\mu(\Delta_2)\big].
\end{equation}
\paragraph{Standard term.} $\mathbb{E}[m_1+m_2]=B\big(x+\tfrac{B}{2}-b\big)$.
\paragraph{Money dimension.} The money outcome is $-b$ with probability $B$ and $0$ otherwise. Comparing against an i.i.d.\ reference,
\begin{equation}
\mathbb{E}[\mu(\Delta_{\text{money}})]
= \eta\,b\,B(1-B) + \eta\lambda(-b)B(1-B)
= -\Lambda\,B(1-B)\,b .
\end{equation}
\paragraph{Good dimension.} Let $G$ denote the realized good (value) and $G'$ its reference copy. Since $G\stackrel{d}{=}G'$,
\begin{equation}
\mathbb{E}[\mu(G-G')] = -\eta(\lambda-1)\,\mathbb{E}[(G-G')^{+}] = -\Lambda\,\mathbb{E}[(G-G')^{+}].
\end{equation}
The good equals $0$ with probability $1-B$ and $x+Y$ ($Y\sim U[0,B]$) with probability $B$. Decomposing,
\begin{equation}
\mathbb{E}[(G-G')^{+}]
= \underbrace{B(1-B)\big(x+\tfrac{B}{2}\big)}_{\text{win vs.\ ref.\ lose}}
+ \underbrace{B^{2}\,\mathbb{E}[(Y-Y')^{+}]}_{\text{both win}}
= B(1-B)\big(x+\tfrac{B}{2}\big) + \frac{B^{3}}{6},
\end{equation}
using $\mathbb{E}[(Y-Y')^{+}]=\tfrac{B}{6}$ for $Y,Y'\sim U[0,B]$ i.i.d.
\paragraph{Aggregate objective.}
\begin{equation}
\boxed{\;
U(B;x)=B\Big(x+\tfrac{B}{2}-b\Big)
-\Lambda\,B(1-B)\,b
-\Lambda\Big[\,B(1-B)\big(x+\tfrac{B}{2}\big)+\tfrac{B^{3}}{6}\,\Big]
\;}
\end{equation}
The last term is the \emph{value-dispersion penalty}: loss aversion over the
uncertainty in the common value \emph{among winners}. It is absent in the
private-value model, where $V=v$ is known.
\subsection{Equilibrium}
\begin{proposition}[First-price wallet game under KR preferences]
There is a unique symmetric, strictly increasing PPE equilibrium
\begin{equation}
\boxed{\ \beta^{*}(x)=\frac{(1-\Lambda)x+\Lambda x^{2}}{1+\Lambda(1-x)}\ }.
\end{equation}
\end{proposition}
\begin{proof}
Substitute $b=\beta(B)$ into $U(B;x)$ and impose the first-order condition
$\partial U/\partial B = 0$ at $B=x$. One obtains the linear ODE
\begin{equation}
x\big[1+\Lambda(1-x)\big]\,\beta'(x)
+\big[1+\Lambda(1-2x)\big]\,\beta(x)
= 2x(1-\Lambda)+3\Lambda x^{2}.
\end{equation}
The left-hand side is a total derivative:
\begin{equation}
\frac{d}{dx}\Big[\,\beta(x)\,x\big(1+\Lambda(1-x)\big)\Big]
= 2x(1-\Lambda)+3\Lambda x^{2}.
\end{equation}
Integrating from $0$ with the boundary condition $\beta(0)=0$,
\begin{equation}
\beta(x)\,x\big(1+\Lambda(1-x)\big)=(1-\Lambda)x^{2}+\Lambda x^{3},
\end{equation}
which rearranges to the stated expression.
\end{proof}
\subsection{Properties}
\begin{itemize}
\item \textbf{Risk-neutral benchmark} $(\Lambda=0)$: $\beta^{*}(x)=x$, recovering $\beta_0$.
\item \textbf{Uniform underbidding.} For all $x\in(0,1)$,
\begin{equation}
\beta^{*}(x)-x=\frac{2\Lambda\,x(x-1)}{1+\Lambda(1-x)}<0,
\end{equation}
so loss aversion \emph{exacerbates the winner's curse}: bids are shaded below the risk-neutral level everywhere.
\item \textbf{Comparative statics.} Writing $\beta^{*}(x)=x\,\dfrac{1-\Lambda(1-x)}{1+\Lambda(1-x)}$, we have
\begin{equation}
\frac{\partial \beta^{*}(x)}{\partial \Lambda}<0.
\end{equation}
\item \textbf{Boundaries.} $\beta^{*}(0)=0$ and $\beta^{*}(1)=1$. Unlike the private-value model (where $\beta^{*}(1)=\tfrac12+\tfrac{\Lambda}{6}>\tfrac12$), there is \emph{no} top-end overbidding.
\item \textbf{Private-value limit.} If $V$ is known, the within-winner dispersion term $B^{3}/6$ vanishes and the problem reduces to the structure of Section~3.
\item \textbf{Source of the distortion.} The distortion is driven entirely by the \emph{multi-dimensional} (good vs.\ money) separation. Under a single, net-surplus reference (below) bidding is undistorted.
\end{itemize}
\subsection{Extensions}
\paragraph{Net-surplus reference (single dimension).}
If preferences are defined over net surplus $\pi=V-b$ with a self-consistent reference,
\begin{equation}
U(B;x)=\Big(x+\tfrac{B}{2}-b\Big)B\big[1-\Lambda(1-B)\big]-\tfrac{\Lambda B^{3}}{6}.
\end{equation}
The first-order condition again reduces to a total derivative and yields
\begin{equation}
\beta^{*}(x)=x,
\end{equation}
so one-dimensional reference dependence does \emph{not} distort the first-price wallet game. The entire effect in Proposition~1 is therefore attributable to the dimensional decomposition of the common value.
\paragraph{General signal distribution.}
Let $X_i\sim F$ (density $f$) on $[0,1]$, and set $P=F(B)$, $I(B)=\int_0^B y\,f(y)\,dy$, and $D(B)=\mathbb{E}\big[(Y-Y')^{+}\mid Y,Y'\le B\big]$. Then
\begin{equation}
U(B;x)=P\!\left(x+\tfrac{I(B)}{P}-b\right)
-\Lambda P(1-P)\,b
-\Lambda\Big[(1-P)\big(xP+I(B)\big)+P^{2}D(B)\Big],
\end{equation}
and the equilibrium function solves the ODE obtained from $\partial U/\partial B=0$ at $B=x$. For $F=U[0,1]$ ($P=B$, $I(B)=\tfrac{B^{2}}{2}$, $D(B)=\tfrac{B}{6}$) this is Proposition~1.
\paragraph{Scaling and $n$ bidders.}
For $V=\alpha(X_1+X_2)$ the problem is homogeneous in $(\alpha,b)$ and
\begin{equation}
\beta^{*}(x)=\alpha\,\frac{(1-\Lambda)x+\Lambda x^{2}}{1+\Lambda(1-x)}.
\end{equation}
With $n$ bidders and $V=\sum_{i=1}^{n}X_i$, the win probability becomes $B^{\,n-1}$ and the dispersion term scales as $B^{\,n-1}\!\big(1-B^{\,n-1}\big)$ and $B^{\,2(n-1)}\mathbb{E}[(Y-Y')^{+}]$; the resulting ODE is solved by the same integrating-factor method.
\paragraph{Second-price / English wallet game (sketch).}
Here the payment equals the highest rival bid, $b(Y)$, and is random conditional on winning. The money dimension then contributes a term
\begin{equation}
-\Lambda\Big[\,B(1-B)\,\mathbb{E}\,b(Y)+B^{2}\,\mathbb{E}\big[(b(Y')-b(Y))^{+}\big]\Big],
\qquad Y,Y'\sim U[0,B],
\end{equation}
which contains $\beta$ inside the gain--loss integral. The equilibrium condition is therefore an integro-differential equation, and the risk-neutral truth-revealing benchmark $\beta_0(x)=2x$ fails---the exact analogue of the Lange--Ratan (2010) ``within-dimension price variance'' result for second-price auctions.



\section{Alternative Explanations for Overbidding}
\label{sec:overbidding}
Laboratory subjects overbid relative to the risk-neutral Bayesian--Nash benchmark
in both private-value (PV) and common-value (CV) auctions, and in both first-price
(FPA) and second-price (SPA) formats. In PV-SPA truthful bidding is a dominant
strategy, so overbidding there is especially revealing. Reference-dependent
preferences of the K\H{o}szegi--Rabin type provide only a setting-dependent and
directionally ambiguous account (approximately bidding up in high-value PV-FPA,
but down in the first-price wallet game); the literature offers several
complementary explanations.
\paragraph{1. Risk preferences and nonlinear probability weighting.}
\emph{Risk aversion} raises bids in PV-FPA \cite{cox1982,coxSmithWalker1988} but is
irrelevant in PV-SPA (truth-telling remains dominant). \emph{Probability
weighting} (overweighting the chance of winning, underweighting payment) inflates
bids in all formats \cite{kahnemanTversky1979}; for auction-specific evidence on
misperceived winning probabilities see \cite{armantierTreich2006}. Risk-seeking or
gambling preferences can likewise raise bids.
\paragraph{2. Non-monetary utility of winning.}
A ``joy of winning'' or consumption value attached to acquiring the object raises
the effective valuation in PV-FPA, PV-SPA, and CV \cite{delgado2008}. Related
mechanisms are competitive arousal (``auction fever'') \cite{kuMalhotraMurnighan2005},
escalation of commitment \cite{rosato2023}, and spiteful/competitive bidding that
values making rivals lose \cite{morgan2003}.
\paragraph{3. Anticipated regret.}
\emph{Loser regret}---regret from having bid too little and lost---shifts optimal
bids upward in FPA and especially SPA \cite{engelbrechtWiggans1989,filizOzbay2007}.
This differs from reference dependence: it compares the chosen action to a
counterfactual, not consumption to an expectations-based reference lottery.
Winner regret can partially offset, but empirically loser regret dominates.
\paragraph{4. Bounded rationality and noisy choice.}
\emph{Quantal response equilibrium} with a flat payoff surface generates
systematic, not zero-mean, overbidding in PV-FPA \cite{goeree2002}.
\emph{Level-$k$/cognitive-hierarchy} reasoning---level-1 agents best-responding to a
naive uniform belief---explains SPA overbidding and part of CV overbidding
\cite{crawfordIriberri2007}. \emph{Learning-direction theory} predicts adjustment
in the profitable direction with systematic overshooting
\cite{seltenBuchta1999}. A sizable share of SPA overbidding appears to be plain
miscalculation, correlated with experience and payoffs \cite{cooperFang2008}.
\paragraph{5. Common-value-specific mechanisms.}
Beyond the classic failure to condition on winning, \emph{cursedness / cursed
equilibrium} formalizes best-responding to the marginal (rather than conditional)
distribution of rivals' signals \cite{eysterRabin2005,kagelLevin1986}.
\emph{Correlation neglect}---failing to account for the correlation between signals
and the event of winning---amplifies the curse \cite{eysterWeizsacker2011}.
Anchoring on one's own signal, and overconfidence in signal precision, further
inflate CV bids.
\paragraph{6. Design and methodological artifacts.}
\emph{Flat payoff surfaces} make over- and underbidding nearly cost-equivalent, so
right-skewed errors or reference normalization surface as average overbidding.
Real-monetary incentives, experience, and cash effects reduce but rarely eliminate
overbidding.
\paragraph{Summary.}
\begin{center}
\begin{tabular}{lccc}
\hline
Explanation & PV-FPA & PV-SPA & CV \\
\hline
Risk aversion                       & yes & no & weak \\
Probability weighting               & yes & yes & yes \\
Joy of winning                      & yes & yes & yes \\
Competitive arousal / spite         & yes & yes & yes \\
Anticipated (loser) regret          & yes & strong & yes \\
QRE / noise                         & strong & yes & yes \\
Level-$k$                           & yes & strong & yes \\
Cursedness / correlation neglect    & no & no & core \\
Learning overshoot                  & yes & yes & strong \\
\hline
\end{tabular}
\end{center}






\section{Introduction}
A robust experimental regularity is that bidders bid above the risk-neutral
Bayesian--Nash benchmark across a wide range of auction environments: first- and
second-price private-value auctions \citep{cox1982,coxSmithWalker1988,goeree2002,
crawfordIriberri2007}, and first- and second-price common-value auctions
\citep{kagelLevin1986,eysterRabin2005,crawfordIriberri2007}. Explanations include
risk aversion \citep{cox1982,coxSmithWalker1988}, anticipated regret
\citep{engelbrechtWiggans1989,filizOzbay2007}, quantal response and level-$k$
reasoning \citep{goeree2002,crawfordIriberri2007}, cursedness \citep{eysterRabin2005},
and a non-monetary utility of winning \citep{delgado2008}.
Expectations-based reference dependence \citep{koszegiRabin2006,koszegiRabin2007},
as applied to sealed-bid auctions by \citet{langeRatan2010}, is a leading candidate.
Its canonical auction formulation, however, is \emph{multi-dimensional}: the good
and money are treated as separate consumption dimensions, which is what generates
the ``attachment'' to the object and the resulting high-value overbidding. This is a
substantial assumption in induced-value laboratory auctions, where the object being
sold is a token whose entire value is a monetary redemption value. There, the good
and money dimensions are economically the same, and it is arguably not legitimate to
evaluate gain--loss separately on each.
This paper takes the induced-value setting seriously and restricts attention to a
\emph{single monetary dimension}, with preferences defined over the net monetary
payoff and a self-consistent (personal-equilibrium) reference point. In a single
dimension, K\H{o}szegi--Rabin preferences reduce to \emph{first-order risk aversion}
\citep{koszegiRabin2007}: the kink in the gain--loss function makes agents averse to
mean-preserving spreads in a way that ordinary expected-utility risk aversion is
not. This channel has sharply different implications from the two-dimensional
model, so the choice of dimensionality is not innocuous.
To this baseline we add a single reference-independent parameter: a direct utility
bonus $w\ge0$ obtained upon winning, the ``joy of winning'' \citep{delgado2008}.
Because the bonus is direct rather than a change in the good's intrinsic value, it
is not itself subject to gain--loss comparison and is well defined whether or not
the good is a separate dimension. The model nests pure loss aversion at $w=0$ and
admits closed-form equilibria in the two canonical first-price settings.
Our results are as follows. In the private-value first-price auction, loss aversion
alone generates overbidding for \emph{every} valuation; the joy-of-winning bonus
amplifies the overbidding and yields a positive bid at a zero valuation, which pure
loss aversion cannot. In the first-price wallet game of \citet{klemperer1998}, loss
aversion alone has \emph{no} effect: the equilibrium coincides with the risk-neutral
benchmark $\beta_0(x)=x$. Any strictly positive joy-of-winning bonus, however,
produces uniform overbidding. This delivers a clean, testable asymmetry:
\emph{private-value} overbidding is consistent with pure loss aversion, while
\emph{common-value} overbidding beyond the winner's curse requires a
joy-of-winning motive. Together with the format-invariant intercept $w/(1-\Lambda)$,
this provides a small set of moments to separate the two motives.
Section~\ref{sec:framework} presents the framework and discusses dimensionality.
Section~\ref{sec:pv} and Section~\ref{sec:wallet} analyze the private-value auction
and the wallet game in the first-price format; Section~\ref{sec:spa} analyzes the two
second-price environments. Section~\ref{sec:id} discusses identification and revenue.
Section~\ref{sec:ext} sketches extensions. An appendix states the two-dimensional
results for comparison.
\section{Framework}
\label{sec:framework}
\subsection{Preferences over a single monetary dimension}
Consumption is the net monetary payoff $\pi$. Following \citet{koszegiRabin2006},
utility relative to reference payoff $r$ is
\begin{equation}
u(\pi\mid r)=\pi+\mu(\pi-r),
\end{equation}
where $\mu$ is piecewise linear,
\begin{equation}
\mu(x)=\begin{cases}\eta\,x & x\ge0,\\ \eta\lambda\,x & x<0,\end{cases}
\qquad \lambda>1,\ \eta>0 .
\end{equation}
We set $\eta=1$ and write $\Lambda\equiv\eta(\lambda-1)>0$. The reference point is the
distribution of consumption induced by the agent's own planned action (a personal
equilibrium), and utility is the expectation of $u(\pi\mid r)$ over an independent
draw $r$ from that distribution. Because $\mu$ is kinked at zero, reference
dependence over a single dimension induces first-order risk aversion
\citep{koszegiRabin2007}, which behaves like risk aversion in first-price auctions
but is a distinct, reference-dependent primitive.
\subsection{Dimensionality and induced values}
\label{sec:dim}
The standard auction application of these preferences
\citep{langeRatan2010} separates the good ($k=1$) and money ($k=2$):
$u(c\mid r)=\sum_k[m_k(c_k)+\mu(m_k(c_k)-m_k(r_k))]$. The good dimension contributes
gain--loss when the bidder does not obtain the object, generating attachment. This
requires $m_1$ to be a consumption dimension psychologically distinct from money.
In an induced-value auction the object's value is a cash redemption value, so
$m_1(v)=v$ is money and the two dimensions coincide. We therefore use the
single-dimension formulation above throughout. A brief two-dimensional comparison is
provided in Appendix~\ref{app:two}. Which dimensionality applies is ultimately an
empirical question; the two models make distinct, testable predictions (see
Section~\ref{sec:id}).
\subsection{Joy of winning}
A bidder who wins receives a direct utility bonus $w\ge0$, modeled additively as
$w\cdot P$, where $P$ is the probability of winning. Being direct and
reference-independent, this specification is valid under either dimensionality.
\begin{remark}[Alternative specification]
If instead $w$ entered the intrinsic value of the good and hence the reference
point, gain--loss would partially offset it. All qualitative conclusions below are
unchanged, with the joy-of-winning effect dampened; we therefore adopt the direct
specification.
\end{remark}
\subsection{Payment rules}
The payment rule determines whether the bidder's payoff conditional on winning depends
on his own action. Under a \emph{first-price} rule the winner pays his own bid, so the
payoff $v-b$ is chosen directly; under a \emph{second-price} rule he pays the rival's
bid and the price is random conditional on winning. We treat the first-price
environments in Sections~\ref{sec:pv}--\ref{sec:wallet} and the second-price
environments in Section~\ref{sec:spa}.
\section{Private-Value First-Price Auction}
\label{sec:pv}
Two symmetric bidders have i.i.d.\ private values $v\sim U[0,1]$. In a symmetric,
strictly increasing equilibrium a bid $b=\beta(P)$ wins with probability $P$. A
bidder with value $v$ who plans such a bid obtains the net payoff $\pi=v-b$ with
probability $P$ and $\pi=0$ otherwise. The reference is an independent draw from
that same two-point distribution. For $v>b$, the gain--loss term is
\begin{equation}
\mathbb{E}[\mu(\pi-\pi')]=-\Lambda\,P(1-P)(v-b),
\end{equation}
so that, adding the joy of winning,
\begin{equation}
U(P;v)=P(v-b)\big[1-\Lambda(1-P)\big]+wP .
\label{eq:pvobj}
\end{equation}
\begin{proposition}[Private-value first-price equilibrium, single dimension]
\label{prop:pv}
The unique symmetric, strictly increasing personal-equilibrium bidding strategy is
\begin{equation}
\boxed{\ \beta^{*}(v)=\frac{w+\tfrac{1-\Lambda}{2}\,v+\tfrac{2\Lambda}{3}\,v^{2}}
{1-\Lambda+\Lambda v}\ }.
\end{equation}
\end{proposition}
\begin{proof}
Substitute $b=\beta(P)$ in \eqref{eq:pvobj}. The first-order condition
$\partial U/\partial P=0$ at $P=v$ gives
\begin{equation}
v(1-\Lambda+\Lambda v)\,\beta'(v)+\big[1-\Lambda+2\Lambda v\big]\beta(v)
=(1-\Lambda)v+2\Lambda v^{2}+w .
\end{equation}
The left-hand side is $\frac{d}{dv}\big[\beta(v)\,v(1-\Lambda+\Lambda v)\big]$.
Integrating and imposing that $\beta(v)\,v(1-\Lambda+\Lambda v)\to0$ as $v\to0$,
\begin{equation}
\beta(v)\,v(1-\Lambda+\Lambda v)=\tfrac{1-\Lambda}{2}v^{2}+\tfrac{2\Lambda}{3}v^{3}+wv,
\end{equation}
which rearranges to the stated expression.
\end{proof}
\begin{corollary}[Overbidding and the low-value intercept]
\label{cor:pvover}
Relative to the risk-neutral benchmark $\tfrac12 v$,
\begin{equation}
\beta^{*}(v)-\tfrac12 v=\frac{w+\tfrac{\Lambda}{6}v^{2}}{1-\Lambda+\Lambda v}\ \ge\ 0 ,
\qquad
\beta^{*}(0)=\frac{w}{1-\Lambda},\qquad \beta^{*}(1)=\tfrac12+\tfrac{\Lambda}{6}+w .
\end{equation}
For $w=0$, loss aversion alone already produces overbidding at every $v>0$.
\end{corollary}
\begin{proof}
Immediate from Proposition~\ref{prop:pv}.
\end{proof}
\begin{remark}[Risk-neutral benchmark]
For $\Lambda=0$, $\beta^{*}(v)=\tfrac12 v+w$, recovering the standard first-price
benchmark $\tfrac12 v$ when $w=0$ and full pass-through of the joy-of-winning bonus.
For $\Lambda>0$ the denominator $1-\Lambda+\Lambda v$ is less than one for $v<1$, so
both channels are amplified relative to the linear case.
\end{remark}
\section{The First-Price Wallet Game}
\label{sec:wallet}
Two symmetric bidders receive private signals $X_i\sim U[0,1]$ i.i.d.\ and the
common value is $V=X_1+X_2$ \citep{klemperer1998,milgromWeber1982}. In a symmetric,
strictly increasing equilibrium a bid $b=\beta(P)$ wins with probability $P$, and
conditional on winning the rival's signal is $Y\sim U[0,P]$. The net payoff is
$\pi=x+Y-b$ with probability $P$ and $0$ otherwise. Compared with
Section~\ref{sec:pv}, the common value adds within-winner uncertainty. Its
gain--loss contribution is
\begin{equation}
\mathbb{E}[\mu(\pi-\pi')]=-\Lambda\Big[P(1-P)\Big(x+\tfrac{P}{2}-b\Big)+\tfrac{P^{3}}{6}\Big],
\end{equation}
using $\mathbb{E}[(Y-Y')^{+}]=P/6$ for $Y,Y'\sim U[0,P]$ i.i.d. Adding the joy of
winning,
\begin{equation}
U(P;x)=P\Big(x+\tfrac{P}{2}-b\Big)\big[1-\Lambda(1-P)\big]-\tfrac{\Lambda P^{3}}{6}+wP .
\label{eq:walletobj}
\end{equation}
\begin{proposition}[Wallet-game equilibrium, single dimension]
\label{prop:wallet}
The unique symmetric, strictly increasing personal-equilibrium bidding strategy is
\begin{equation}
\boxed{\ \beta^{*}(x)=\frac{w+(1-\Lambda)x+\Lambda x^{2}}{1-\Lambda+\Lambda x}\ }.
\end{equation}
\end{proposition}
\begin{proof}
Substitute $b=\beta(P)$ in \eqref{eq:walletobj}. The first-order condition at $P=x$
gives
\begin{equation}
x(1-\Lambda+\Lambda x)\,\beta'(x)+\big[1-\Lambda+2\Lambda x\big]\beta(x)
=2x(1-\Lambda)+3\Lambda x^{2}+w .
\end{equation}
The left-hand side equals
$\frac{d}{dx}\big[\beta(x)\,x(1-\Lambda+\Lambda x)\big]$. Integrating with
$\beta(x)\,x(1-\Lambda+\Lambda x)\to0$ as $x\to0$,
\begin{equation}
\beta(x)\,x(1-\Lambda+\Lambda x)=(1-\Lambda)x^{2}+\Lambda x^{3}+wx,
\end{equation}
which rearranges to the stated expression.
\end{proof}
\begin{corollary}[No loss-aversion distortion; joy of winning overbids]
\label{cor:wallet}
For all $x\in[0,1]$,
\begin{equation}
\beta^{*}(x)-x=\frac{w}{1-\Lambda+\Lambda x}\ \ge\ 0 ,
\qquad
\beta^{*}(0)=\frac{w}{1-\Lambda},\qquad \beta^{*}(1)=1+w .
\end{equation}
In particular, if $w=0$ then $\beta^{*}(x)=x$ exactly: with a single monetary
dimension, loss aversion does \emph{not} distort the wallet-game equilibrium. If
$w>0$, the bid exceeds $x$ uniformly.
\end{corollary}
\begin{proof}
Immediate from Proposition~\ref{prop:wallet}.
\end{proof}
\begin{remark}[Interpretation]
The winner's curse makes the common value random conditional on winning; the
gain--loss terms in \eqref{eq:walletobj} then cancel at the margin and the
risk-neutral bid $\beta_0(x)=x$ survives. Thus loss aversion can rationalize
private-value overbidding but not, by itself, common-value overbidding. The latter
requires the joy-of-winning motive (or a non-reference mechanism such as cursedness
or bounded rationality).
\end{remark}
\section{Second-Price Auctions}
\label{sec:spa}
Under a second-price rule the winner pays the rival's bid rather than his own, so the
payment is random conditional on winning. This \emph{within-dimension price variance}
is the feature that distinguishes the second-price from the first-price format. We
analyze the two canonical environments of the preceding sections---the private-value
auction and the wallet game---under the same single-monetary-dimension preferences.
\subsection{Independent-Private-Value Second-Price Auction}
\label{sec:ipvspa}
Two symmetric bidders have i.i.d.\ private values $v\sim U[0,1]$. In a symmetric,
strictly increasing equilibrium $\beta$, a bidder who plans a bid that wins with
probability $P$ (i.e.\ $b=\beta(P)$) wins when the rival's value $Y<P$ and pays the
rival's bid $\beta(Y)$, with $Y\sim U[0,P]$. His net payoff is therefore
\begin{equation}
\pi=v-\beta(Y)\ \text{with probability }P,\quad Y\sim U[0,P];
\qquad \pi=0\ \text{with probability }1-P,
\end{equation}
and---as in Section~\ref{sec:pv}---the reference point is an independent draw from this
same distribution. Since outcomes and reference are i.i.d.,
$\mathbb{E}[\mu(\pi-\pi')]=-\Lambda\,\mathbb{E}[(\pi'-\pi)^+]$. Decomposing the positive
difference into the one-wins/one-loses and both-win events, and using
$\iint(\beta(y)-\beta(y'))^{+}dy\,dy'=\int_0^P(2y-P)\beta(y)\,dy$ for increasing
$\beta$, the interim objective is
\begin{equation}
U(P;v)=vP-\bar\beta(P)
-\Lambda\Big[(1-P)\!\int_0^P\!|v-\beta(y)|\,dy+\int_0^P(2y-P)\beta(y)\,dy\Big]+wP,
\label{eq:spaobj}
\end{equation}
where $\bar\beta(P)\equiv\int_0^P\beta(y)\,dy$.
\begin{proposition}[No loss-aversion distortion in the private-value SPA]
\label{prop:spa0}
Let $0<\Lambda<1$ and $w=0$. The unique symmetric, strictly increasing
personal-equilibrium bidding strategy is
\begin{equation}
\boxed{\ \beta^{*}(v)=v\ }.
\end{equation}
\end{proposition}
\begin{proof}
At $w=0$ the equilibrium does not overbid, so $|v-\beta(y)|=v-\beta(y)$ for $y\le v$.
Differentiating \eqref{eq:spaobj} at $P=v$, and using
$\partial_P\int_0^P(2y-P)\beta(y)\,dy=P\beta(P)-\bar\beta(P)$,
\begin{equation}
\frac{\partial U}{\partial P}\Big|_{P=v}
=v-\beta(v)-\Lambda\Big[-\!\int_0^v\!\big(v-\beta(y)\big)dy+(1-v)\big(v-\beta(v)\big)
+v\beta(v)-\bar\beta(v)\Big].
\end{equation}
The term in brackets equals $v-2v^{2}+(2v-1)\beta(v)$, so the first-order condition
becomes
\begin{equation}
\big(v-\beta(v)\big)\big(1-\Lambda+2\Lambda v\big)=0 .
\end{equation}
Since $1-\Lambda+2\Lambda v\ge 1-\Lambda>0$, the unique solution is $\beta(v)=v$.
Moreover $U''(v)=-(1-\Lambda+2\Lambda v)<0$ and $\partial_P U<0$ for $P>v$, so $P=v$
is the global maximizer; truthful bidding is the unique symmetric equilibrium.
\end{proof}
\begin{proposition}[Joy of winning in the private-value SPA]
\label{prop:spa}
(i) If $\Lambda=0$, the equilibrium is $\beta^{*}(v)=v+w$ exactly. (ii) Let
$0<\Lambda<1$ and let $w>0$ be small. Then
\begin{equation}
\beta^{*}(v)=v+\frac{w}{1+\Lambda}+O(w^{2}),
\end{equation}
so bids exceed the risk-neutral benchmark at every $v$, and loss aversion
\emph{dampens} the joy-of-winning premium (the denominator $1+\Lambda$ exceeds one,
whereas in the private-value FPA the corresponding denominator $1-\Lambda+\Lambda v$
is below one).
\end{proposition}
\begin{proof}
(i) For $\Lambda=0$, \eqref{eq:spaobj} reduces to $U(P;v)=vP-\bar\beta(P)+wP$, whose
first-order condition at $P=v$ is $v-\beta(v)+w=0$, giving $\beta(v)=v+w$.

(ii) The exact first-order condition at the symmetric equilibrium, retaining the kink,
is
\begin{equation}
v-\beta(v)-\Lambda\Big[-\!\int_0^v\!|v-\beta(y)|\,dy+(1-v)\big|\beta(v)-v\big|
+v\beta(v)-\bar\beta(v)\Big]+w=0 .
\label{eq:spaexactfoc}
\end{equation}
Write $\beta(y)=y+w\phi(y)+O(w^{2})$. Since $\beta(v)-v=O(w)$, the region on which
$|v-\beta(y)|\neq v-\beta(y)$ has width $O(w)$ and contributes only $O(w^{2})$, so
$\int_0^v|v-\beta(y)|dy=\tfrac{v^{2}}{2}-w\int_0^v\phi(y)\,dy+O(w^{2})$. With
$v\beta(v)=v^{2}+wv\phi(v)+O(w^{2})$, $\bar\beta(v)=\tfrac{v^{2}}{2}+w\int_0^v\phi+O(w^{2})$
and $|\beta(v)-v|=w\phi(v)$, the bracket in \eqref{eq:spaexactfoc} equals
$w\phi(v)+O(w^{2})$. Substituting, \eqref{eq:spaexactfoc} becomes
$-w\phi(v)-\Lambda w\phi(v)+w+O(w^{2})=0$, i.e.\ $\phi(v)=1/(1+\Lambda)+O(w)$; the
second-order condition holds for small $w$ by continuity from part (i).
\end{proof}
\begin{remark}
On the interior $v\ge c$ the exact equilibrium is the constant-offset bid
$\beta^{*}(v)=v+c$, where $c$ solves $\Lambda c^{2}-(1+\Lambda)c+w=0$, i.e.
$c=\frac{(1+\Lambda)-\sqrt{(1+\Lambda)^{2}-4\Lambda w}}{2\Lambda}=\frac{w}{1+\Lambda}+O(w^{2})$.
On the boundary layer $v<c=O(w)$ every win is a net loss, so that characterization
breaks down; we therefore report the global result to leading order in $w$.
\end{remark}
\subsection{The Second-Price Wallet Game}
\label{sec:walletspsa}
Signals are $X_i\sim U[0,1]$ and the value is $V=X_1+X_2$. The risk-neutral
second-price (English) equilibrium is $\beta_0(x)=2x$: a winner at price $P$ earns
$x+\mathbb{E}[Y\mid Y<P]-P$, and bidding $2x$ equates the price to the value expected
conditional on winning.
\begin{proposition}[Wallet-game SPA under reference-dependent preferences]
\label{prop:walletspsa}
The wallet game is isomorphic to the private-value SPA. If the rival uses
$\beta(y)=2y+c$, a bidder with signal $x$ who wins at price $\beta(Y)$ earns
$x+Y-\beta(Y)=x-Y-c$ on $Y\in[0,P]$, which is exactly the private-value objective with
value $x$ and price $Y+c$. Consequently: (i) at $w=0$, $\beta^{*}(x)=2x$; (ii) at
$\Lambda=0$, $\beta^{*}(x)=2x+w$; and (iii) for small $w>0$,
$\beta^{*}(x)=2x+\frac{w}{1+\Lambda}+O(w^{2})$.
\end{proposition}
\begin{proof}
Immediate from the isomorphism and Propositions~\ref{prop:spa0}--\ref{prop:spa}.
\end{proof}
\begin{remark}[Interpretation]
With a single monetary dimension, loss aversion distorts neither second-price format:
the risk-neutral benchmarks $\beta_0(v)=v$ and $\beta_0(x)=2x$ survive at $w=0$. This is
the second-price counterpart of the first-price wallet-game result
(Corollary~\ref{cor:wallet}) and stands in sharp contrast to the private-value FPA,
where loss aversion alone overbids (Corollary~\ref{cor:pvover}). Overbidding in either
second-price format therefore requires the joy-of-winning motive, at a rate damped by
loss aversion.
\end{remark}
\section{Identification and Revenue}
\label{sec:id}
The single-dimension model yields the following testable implications.
\begin{itemize}
\item \textbf{Private-value overbidding under pure loss aversion.} By
Corollary~\ref{cor:pvover}, the first-price bid satisfies $\beta^*(v)>\tfrac12 v$ for
all $v>0$ even when $w=0$. The local slope of the bid function identifies $\Lambda$.
\item \textbf{Common-value overbidding requires a joy of winning.} By
Corollary~\ref{cor:wallet}, the first-price wallet-game bid equals $x$ exactly when
$w=0$ and exceeds $x$ uniformly when $w>0$. Hence any \emph{excess} common-value
overbidding (beyond the winner's curse) identifies $w>0$.
\item \textbf{No second-price distortion from loss aversion.} By
Propositions~\ref{prop:spa0} and \ref{prop:walletspsa}, when $w=0$ the second-price
private-value equilibrium is truthful, $\beta^{*}(v)=v$, and the second-price wallet
game yields $\beta^{*}(x)=2x$; loss aversion alone distorts neither. Overbidding in
either second-price format therefore again identifies $w>0$, and the \emph{contrast}
between a distorted private-value FPA and an undistorted private-value SPA isolates the
dimensionality of the gain--loss evaluation.
\item \textbf{Positive low-value bids.} A strictly positive low-value bid that is
replicated across formats requires $w>0$, which pure loss aversion cannot generate. Its
\emph{level} is format-specific---$\beta^{*}(0)=w/(1-\Lambda)$ in the FPA and, to
leading order, $w/(1+\Lambda)$ in the second-price formats---so the \emph{sign}, not
the level, is the format-invariant identifying moment.
\end{itemize}
These implications also discriminate against the two-dimensional model, whose sign
pattern is different (Appendix~\ref{app:two}).
Expected revenue in the private-value auction is
$R=2\int_0^1 x\,\beta^*(x)\,dx$, which is strictly increasing in $w$; the
wallet-game revenue is defined analogously from the first-order statistic of the two
signals. Because the joy-of-winning bonus shifts bids up while loss aversion
operates differently across formats, cross-format revenue comparisons inherit the
identification asymmetry above.
\section{Extensions}
\label{sec:ext}
\paragraph{Second-price auctions.}
These are analyzed in Section~\ref{sec:spa}. Under a single monetary dimension, loss
aversion alone leaves the risk-neutral benchmark intact---truthful bidding
$\beta_0(v)=v$ in the private-value case and $\beta_0(x)=2x$ in the wallet game---while
the joy-of-winning bonus produces overbidding at the damped rate $w/(1+\Lambda)$. For
general signal distributions the price $\beta(Y)$ enters the gain--loss integral and
the equilibrium condition becomes an equation in $\beta$ and $\bar\beta$ rather than a
closed form.
\paragraph{General signal distributions and $n$ bidders.}
For $X_i\sim F$ on $[0,1]$ and $V=\alpha(X_1+X_2)$, the problem is homogeneous in
$(\alpha,b)$ and the equilibrium scales by $\alpha$. With $n$ bidders and
$V=\sum_{i=1}^{n}X_i$, the win probability becomes $P^{\,n-1}$ and the dispersion
term scales accordingly; the same integrating-factor method applies.
\paragraph{Toward structural estimation.}
Embedding $(w,\Lambda)$ in a structural estimator that also allows for noisy bidding
\citep{goeree2002} or cursedness \citep{eysterRabin2005} would let researchers
quantify the relative contributions of joy of winning, loss aversion, and bounded
rationality. The single-dimension baseline is the appropriate starting point for
induced-value data.
\appendix
\section{Two-Dimensional Comparison}
\label{app:two}
For completeness we record the multi-dimensional (good, money) results of the
private-value first-price auction and the wallet game, which require the good to be
a separate consumption dimension.
\begin{itemize}
\item \emph{Private-value FPA:}
\begin{equation}
\beta^{*}_{2}(v)=\frac{w+\tfrac{1-\Lambda}{2}v+\tfrac{2\Lambda}{3}v^{2}}{1+\Lambda(1-v)},
\qquad
\beta^{*}_{2}(v)-\tfrac12 v=\frac{w-\Lambda v+\tfrac{7\Lambda}{6}v^{2}}{1+\Lambda(1-v)} .
\end{equation}
With $w=0$, the equilibrium \emph{underbids} for $v<6/7$ and overbids for $v>6/7$;
global overbidding requires $w\ge3\Lambda/14$.
\item \emph{Wallet game:}
\begin{equation}
\beta^{*}_{2}(x)=\frac{w+(1-\Lambda)x+\Lambda x^{2}}{1+\Lambda(1-x)},
\qquad
\beta^{*}_{2}(x)-x=\frac{w-2\Lambda x(1-x)}{1+\Lambda(1-x)} .
\end{equation}
With $w=0$, the equilibrium \emph{underbids} uniformly; global overbidding requires
$w\ge\Lambda/2$.
\item \emph{Private-value SPA:} against a truthful rival, a bidder with value $v$ who
chooses win probability $P$ obtains
$U(P;v)=vP-\tfrac{P^{2}}{2}-\Lambda P(1-P)v-\Lambda\big(\tfrac{P^{2}}{2}-\tfrac{P^{3}}{3}\big)+wP$,
so $U'(v)=-\Lambda v(2-3v)+w$. At $w=0$ the equilibrium \emph{underbids} for $v<2/3$
and overbids for $v>2/3$; global overbidding requires $w\ge\Lambda/3$.
\item \emph{Wallet SPA:} against the risk-neutral benchmark $\beta_0(x)=2x$,
$U'(x)=-\Lambda x(4-5x)+w$, so at $w=0$ the equilibrium \emph{underbids} for $x<4/5$
and overbids for $x>4/5$; global overbidding requires $w\ge4\Lambda/5$.
\end{itemize}
The denominators ($1+\Lambda(1-v)$, $1+\Lambda(1-x)$) differ from the single-dimension
denominators ($1-\Lambda+\Lambda v$, $1-\Lambda+\Lambda x$), and the sign of the
pure-loss-aversion effect reverses. In the second-price formats the two-dimensional
model likewise produces a mixed sign pattern at $w=0$ rather than the undistorted
single-dimension benchmarks $\beta^{*}(v)=v$ and $\beta^{*}(x)=2x$. The two
specifications can thus be distinguished empirically by the behavior of bids at low
valuations: single-dimension predicts overbidding near $v=0$, two-dimension predicts
underbidding.
\begin{thebibliography}{99}
\bibitem[Banerji and Gupta(2014)]{banerji2014}
Banerji, A., and N. Gupta (2014): ``Detection, Identification, and Estimation of
Loss Aversion: Evidence from an Auction Experiment,'' \emph{American Economic
Journal: Microeconomics}, 6(1), 91--133.
\bibitem[Cox, Roberson, and Smith(1982)]{cox1982}
Cox, J. C., B. Roberson, and V. L. Smith (1982): ``Theory and Behavior of Single
Object Auctions,'' \emph{Research in Experimental Economics}, 2, 1--43.
\bibitem[Cox, Smith, and Walker(1988)]{coxSmithWalker1988}
Cox, J. C., V. L. Smith, and J. M. Walker (1988): ``Theory and Individual Behavior
of First-Price Auctions,'' \emph{Journal of Risk and Uncertainty}, 1(1), 61--99.
\bibitem[Crawford and Iriberri(2007)]{crawfordIriberri2007}
Crawford, V. P., and N. Iriberri (2007): ``Level-$k$ Auctions: Can a Nonequilibrium
Model of Strategic Thinking Explain the Winner's Curse and Overbidding in
Private-Value Auctions?,'' \emph{Econometrica}, 75(6), 1721--1770.
\bibitem[Delgado et al.(2008)]{delgado2008}
Delgado, M. R., A. Schotter, E. Y. Ozbay, and E. A. Phelps (2008): ``Understanding
Overbidding: Using the Neural Circuitry of Reward to Design Economic Auctions,''
\emph{Science}, 321(5897), 1849--1852.
\bibitem[Engelbrecht-Wiggans(1989)]{engelbrechtWiggans1989}
Engelbrecht-Wiggans, R. (1989): ``The Effect of Regret on Optimal Bidding in
Auctions,'' \emph{Management Science}, 35(6), 685--692.
\bibitem[Eyster and Rabin(2005)]{eysterRabin2005}
Eyster, E., and M. Rabin (2005): ``Cursed Equilibrium,'' \emph{Econometrica}, 73(5),
1623--1672.
\bibitem[Filiz-Ozbay and Ozbay(2007)]{filizOzbay2007}
Filiz-Ozbay, E., and E. Y. Ozbay (2007): ``Auctions with Anticipated Regret: Theory
and Experiment,'' \emph{American Economic Review}, 97(4), 1407--1418.
\bibitem[Goeree, Holt, and Palfrey(2002)]{goeree2002}
Goeree, J. K., C. A. Holt, and T. R. Palfrey (2002): ``Quantal Response Equilibrium
and Overbidding in Private-Value Auctions,'' \emph{Journal of Economic Theory},
104(1), 247--272.
\bibitem[Kagel and Levin(1986)]{kagelLevin1986}
Kagel, J. H., and D. Levin (1986): ``The Winner's Curse and Public Information in
Common Value Auctions,'' \emph{American Economic Review}, 76(5), 894--920.
\bibitem[Klemperer(1998)]{klemperer1998}
Klemperer, P. (1998): ``Auctions with Almost Common Values: The `Wallet Game' and
its Applications,'' \emph{European Economic Review}, 42(3--5), 757--769.
\bibitem[K\H{o}szegi and Rabin(2006)]{koszegiRabin2006}
K\H{o}szegi, B., and M. Rabin (2006): ``A Model of Reference-Dependent
Preferences,'' \emph{Quarterly Journal of Economics}, 121(4), 1133--1165.
\bibitem[K\H{o}szegi and Rabin(2007)]{koszegiRabin2007}
K\H{o}szegi, B., and M. Rabin (2007): ``Reference-Dependent Risk Attitudes,''
\emph{American Economic Review}, 97(4), 1047--1073.
\bibitem[Lange and Ratan(2010)]{langeRatan2010}
Lange, A., and A. Ratan (2010): ``Multi-Dimensional Reference-Dependent Preferences
in Sealed-Bid Auctions---How (Most) Laboratory Experiments Differ from the Field,''
\emph{Games and Economic Behavior}, 68(2), 634--645.
\bibitem[Milgrom and Weber(1982)]{milgromWeber1982}
Milgrom, P. R., and R. J. Weber (1982): ``A Theory of Auctions and Competitive
Bidding,'' \emph{Econometrica}, 50(5), 1089--1122.
\end{thebibliography}
\end{document}





